%% SYNTHETIC TEST UNIT generated by writing/tools/make_outline_book.py from FINAL_OUTLINE_v2 -- not book text
\chapter{LaMET 与格点因子化：季向东等人的因子化证明}\label{ch:lamet}
本单元为测试用的合成文本（由 FINAL\_OUTLINE\_v2 自动生成），陈述内容从略。

\section{第 1 节（测试）}\label{sec:lamet-intro}

\begin{conjecture}[LaMET 原理（非微扰形式） \tagC]\label{conj:lamet-principle}
（测试用合成陈述：内容从略。）
\uses{def:pdf-quark, def:pdf-lightray}
\end{conjecture}
\paragraph{地位与缺口}测试：本陈述的地位由\cref{conj:ee-np-ope}、\cref{asm:intro-np-ope}、\cref{rem:power-np-status}、\cref{cor:lamet-principle-pt}、\cref{asm:lamet-uniform}、\cref{asm:lamet-twist4}、\cref{asm:lamet-large-z}、\cref{thm:lamet-hadronic}、\cref{prop:lamet-bvz}、\cref{thm:yzliu-gn-convergence}决定。

\begin{remark}[LaMET 原理的几种表述及其可证明的内容]\label{rem:lamet-formulations}
（测试用合成陈述：内容从略。）
% uses: none
\end{remark}

\begin{proposition}[格点可计算的矩 \tagR]\label{prop:lattice-moments}
（$R_{\mathrm{OPE}}$：以 OPE 为前提而严格。QCD 中含夸克的微扰 OPE 是全阶论证（\cref{thm:dis-ope}），严格证明只对标量理论与纯 Yang--Mills 的规范不变算符给出；对强子态还须非微扰 OPE（\cref{asm:intro-np-ope}）。）
（测试用合成陈述：内容从略。）
非微扰地位：\cref{asm:intro-wightman}、\cref{asm:intro-composite}、\cref{asm:intro-np-ope}。
\uses{thm:pdf-moments-local, prop:dis-dispersion, prop:dis-dispersion-np, thm:dis-moments, thm:dis-moments-np, prop:power-rearrangement}
\end{proposition}

\begin{definition}[准分布、赝分布与格点截面]\label{def:lattice-quasi}
（测试用合成陈述：内容从略。）
\uses{def:prelim-wilson, def:pdf-quark, def:pdf-lightray}
\end{definition}

\begin{proposition}[准关联函数的支撑性质（逐图） \tagR]\label{prop:lamet-support}
（$R_{\mathrm{PT}}$：在重整化微扰论中对红外正规化的部分子态严格。）
（测试用合成陈述：内容从略。）
\uses{def:lattice-quasi}
\end{proposition}

\section{第 2 节（测试）}\label{sec:lamet-principle}

\begin{corollary}[Ioffe 时间分布是指数型整函数 \tagR]\label{cor:lamet-entire}
（$R_{\mathrm{PT}}$：在重整化微扰论中对红外正规化的部分子态严格。）
（测试用合成陈述：内容从略。）
\uses{prop:lamet-support}
\end{corollary}

\begin{example}[自由夸克的准 PDF：树图与单圈非物理区]\label{ex:lattice-quasi-tree}
（测试用合成陈述：内容从略。）
\uses{def:lattice-quasi, prop:prelim-wilson-rules}
\end{example}

\begin{assumption}[含类空辅助静态色源的 QCD 的可重整性 \tagA]\label{asm:lamet-aux-renorm}
（测试用合成陈述：内容从略。）
地位：测试状态句。
文献：测试。
被使用于：\cref{prop:lattice-renormalizability}、\cref{asm:lamet-ope}、\cref{asm:lamet-gap}、\cref{conj:lamet-decay-qcd}、\cref{asm:lamet-staple}、\cref{thm:yzliu-soft-hqet}。
解除途径：测试解除途径。
\uses{def:prelim-qcd, def:prelim-brst, def:prelim-wilson}
\end{assumption}
\begin{proofsketch}测试。\end{proofsketch}

\begin{proposition}[Wilson 线夸克双线性算符的乘法可重整性 \tagA]\label{prop:lattice-renormalizability}
（测试用合成陈述：内容从略。）
（本命题只涉及部分子态。）
\uses{asm:lamet-aux-renorm, def:prelim-wilson, def:prelim-msbar, thm:prelim-composite-renorm}
\end{proposition}
\begin{proofsketch}测试。\end{proofsketch}
\paragraph{地位与缺口}测试：本陈述的地位由\cref{asm:lamet-aux-renorm}决定。

\begin{proposition}[类空与类时 Wilson 线算符的反常维度相同 \tagA]\label{prop:lamet-ad-spacelike}
（测试用合成陈述：内容从略。）
\uses{prop:lattice-renormalizability}
\end{proposition}
\begin{proofsketch}测试。\end{proofsketch}
\paragraph{地位与缺口}测试：本陈述的地位由\cref{asm:lamet-aux-renorm}、\cref{prop:lattice-renormalizability}决定。

\section{第 3 节（测试）}\label{sec:lattice-moments}

\begin{definition}[准关联函数的重整化方案]\label{def:lamet-schemes}
（测试用合成陈述：内容从略。）
\uses{def:lattice-quasi, prop:lattice-renormalizability, def:prelim-msbar}
\end{definition}

\begin{lemma}[残余线性项 $e^{-m_0|z|}$ 的 Lorentz 抹平]\label{lem:lamet-lorentzian}
（测试用合成陈述：内容从略。）
\uses{def:lattice-quasi}
\end{lemma}

\begin{proposition}[Wilson 线自能的线性重整子与领头重整子重求和 \tagP]\label{prop:lamet-lrr}
（测试用合成陈述：内容从略。）
\uses{lem:lamet-lorentzian, def:lamet-schemes, def:power-borel, prop:power-renormalon-examples, prop:lattice-renormalizability}
\end{proposition}
\paragraph{地位与缺口}测试：见文献。

\begin{proposition}[欧氏格点矩阵元与 Minkowski 矩阵元的等同 \tagA]\label{prop:lamet-euclid-minkowski}
（测试用合成陈述：内容从略。）
\uses{def:lattice-quasi}
\end{proposition}
\begin{proofsketch}测试。\end{proofsketch}
\paragraph{地位与缺口}测试：见文献。

\begin{proposition}[准 PDF 匹配的单圈结果 \tagR]\label{prop:lamet-oneloop}
（$R_{\mathrm{PT}}$：在重整化微扰论中对红外正规化的部分子态严格。）
（测试用合成陈述：内容从略。）
\uses{def:lattice-quasi, def:lamet-schemes, def:prelim-msbar, prop:pdf-partonic-dr, prop:pdf-lo-kernels, def:prelim-plus}
\end{proposition}

\section{第 4 节（测试）}\label{sec:lattice-quasi}

\begin{proposition}[单圈准 PDF 的尾部、矩与卷积定义域 \tagR]\label{prop:lamet-tails}
（$R_{\mathrm{PT}}$：在重整化微扰论中对红外正规化的部分子态严格。）
（测试用合成陈述：内容从略。）
\uses{prop:lamet-oneloop}
\end{proposition}

\begin{proposition}[非单态准 PDF 匹配核的两圈结果 \tagR]\label{prop:lamet-twoloop}
（$R_{\mathrm{PT}}$：在重整化微扰论中对红外正规化的部分子态严格。）
（测试用合成陈述：内容从略。）
\uses{prop:lamet-oneloop}
\end{proposition}

\begin{proposition}[极化、横度与扭度 3 准分布的匹配 \tagP]\label{prop:lamet-polarized}
（测试用合成陈述：内容从略。）
\uses{def:lattice-quasi, def:pdf-polarized, prop:lattice-renormalizability}
\end{proposition}
\paragraph{地位与缺口}测试：本陈述的地位由\cref{def:prelim-gamma5}决定。

\begin{assumption}[类空 Wilson 线算符的短距离 OPE（重整化微扰论） \tagA]\label{asm:lamet-ope}
（测试用合成陈述：内容从略。）
地位：测试状态句。
文献：测试。
被使用于：\cref{thm:lamet-coordinate}、\cref{thm:lamet-parton-momentum}、\cref{prop:lamet-tmc}、\cref{prop:lamet-sdf}、\cref{thm:lamet-quasi-gpd}。
解除途径：测试解除途径。
\uses{prop:lattice-renormalizability, asm:lamet-aux-renorm, def:dis-twist}
\end{assumption}
\begin{proofsketch}测试。\end{proofsketch}

\begin{theorem}[坐标空间因子化（赝分布形式） \tagR]\label{thm:lamet-coordinate}
（$R_{\mathrm{cond}}$：在所列编号假设之下证明完整；结论的地位不高于其中最弱的假设。）
前提：\cref{asm:lamet-ope}、\cref{prop:lattice-renormalizability}。
（测试用合成陈述：内容从略。）
（本命题只涉及部分子态。）
\uses{asm:lamet-ope, prop:lattice-renormalizability, prop:lamet-support, cor:lamet-entire, def:lattice-quasi, def:pdf-lightray, thm:pdf-moments-local}
\end{theorem}

\section{第 5 节（测试）}\label{sec:lamet-renorm}

\begin{theorem}[动量空间形式与部分子动量 \tagR]\label{thm:lamet-parton-momentum}
（$R_{\mathrm{cond}}$：在所列编号假设之下证明完整；结论的地位不高于其中最弱的假设。）
前提：\cref{asm:lamet-ope}、\cref{thm:lamet-coordinate}、\cref{prop:lattice-renormalizability}。
（测试用合成陈述：内容从略。）
\uses{thm:lamet-coordinate, asm:lamet-ope, prop:lattice-renormalizability, def:lattice-quasi, prop:lamet-tails}
\end{theorem}

\begin{remark}[以强子动量为变量的早期匹配形式不成立 \tagV]\label{rem:lamet-early-form}
（测试用合成陈述：内容从略。）
\uses{thm:lamet-parton-momentum, prop:lamet-oneloop}
\end{remark}
\paragraph{地位与缺口}测试：见文献。

\begin{proposition}[光锥规范下的共线减除路线（Ma–Qiu） \tagA]\label{prop:lamet-cfp}
（测试用合成陈述：内容从略。）
（本命题只涉及部分子态。）
\uses{def:lattice-quasi, prop:lattice-renormalizability, def:pdf-quark, thm:pdf-dglap}
\end{proposition}
\begin{proofsketch}测试。\end{proofsketch}
\paragraph{地位与缺口}测试：本陈述的地位由\cref{thm:lamet-quasi-tmd}决定。

\begin{theorem}[“好的格点截面”的因子化（Ma–Qiu） \tagA]\label{thm:lamet-lcs}
（测试用合成陈述：内容从略。）
\uses{thm:dis-ope, thm:pdf-moments-local, cor:lamet-entire, def:lattice-quasi}
\end{theorem}
\begin{proofsketch}测试。\end{proofsketch}
\paragraph{地位与缺口}测试：本陈述的地位由\cref{asm:lamet-ope}、\cref{thm:dis-moments}决定。

\begin{theorem}[LaMET 匹配定理（部分子态、领头幂次、全阶） \tagA]\label{thm:lattice-lamet}
（测试用合成陈述：内容从略。）
（本命题只涉及部分子态。）
\uses{thm:lamet-parton-momentum, prop:lamet-cfp, prop:lattice-renormalizability, prop:lamet-tails}
\end{theorem}
\begin{proofsketch}测试。\end{proofsketch}
\paragraph{地位与缺口}测试：本陈述的地位由\cref{asm:lamet-ope}决定。

\section{第 6 节（测试）}\label{sec:lamet-oneloop}

\begin{proposition}[单态夸克与胶子准 PDF：带混合的匹配 \tagA]\label{prop:lamet-gluon-singlet}
（测试用合成陈述：内容从略。）
\uses{thm:lattice-lamet, prop:lattice-renormalizability, prop:lamet-tails, prop:pdf-lo-kernels}
\end{proposition}
\begin{proofsketch}测试。\end{proofsketch}
\paragraph{地位与缺口}测试：见文献。

\begin{proposition}[混合方案的单圈匹配核 \tagR]\label{prop:lamet-hybrid-kernel}
（$R_{\mathrm{PT}}$：在重整化微扰论中对红外正规化的部分子态严格。）
（测试用合成陈述：内容从略。）
\uses{prop:lamet-oneloop, def:lamet-schemes}
\end{proposition}

\begin{example}[混合核的归一化]\label{ex:lamet-hybrid-normalization}
（测试用合成陈述：内容从略。）
\uses{prop:lamet-hybrid-kernel}
\end{example}

\begin{proposition}[动量空间匹配在何种重整化方案中成立 \tagA]\label{prop:lamet-scheme-validity}
（测试用合成陈述：内容从略。）
\uses{thm:lattice-lamet, def:lamet-schemes, prop:lamet-hybrid-kernel}
\end{proposition}
\begin{proofsketch}测试。\end{proofsketch}
\paragraph{地位与缺口}测试：本陈述的地位由\cref{thm:lattice-lamet}决定。

\begin{remark}[比值与 RI/MOM 方案在大 $|z|$ 处的动量空间匹配 \tagC]\label{rem:lamet-ratio-large-z}
（测试用合成陈述：内容从略。）
\uses{prop:lamet-scheme-validity}
\end{remark}
\paragraph{地位与缺口}测试：见文献。

\section{第 7 节（测试）}\label{sec:lattice-matching}

\begin{assumption}[谱隙与边界值正则性 \tagC]\label{asm:lamet-gap}
（测试用合成陈述：内容从略。）
地位：测试状态句。
文献：测试。
被使用于：\cref{thm:lamet-exp-decay}、\cref{conj:lamet-decay-qcd}。
解除途径：测试解除途径。
\uses{asm:lamet-aux-renorm}
\end{assumption}

\begin{theorem}[局域性蕴含类空关联函数的指数衰减（Ji–Liu–Su） \tagR]\label{thm:lamet-exp-decay}
（$R_{\mathrm{cond}}$：在所列编号假设之下证明完整；结论的地位不高于其中最弱的假设。）
前提：\cref{asm:intro-wightman}、\cref{asm:lamet-gap}。
（测试用合成陈述：内容从略。）
\uses{asm:intro-wightman, asm:lamet-gap}
\end{theorem}

\begin{conjecture}[QCD 准分布矩阵元的长程渐近形式 \tagC]\label{conj:lamet-decay-qcd}
（测试用合成陈述：内容从略。）
\uses{thm:lamet-exp-decay, asm:lamet-gap, asm:lamet-aux-renorm, def:lattice-quasi}
\end{conjecture}
\paragraph{地位与缺口}测试：见文献。

\begin{lemma}[准 PDF 的解析带]\label{lem:lamet-strip}
（测试用合成陈述：内容从略。）
\uses{def:lattice-quasi}
\end{lemma}

\begin{proposition}[Fourier 截断误差界]\label{prop:lamet-truncation}
（测试用合成陈述：内容从略。）
\uses{lem:lamet-strip}
\end{proposition}

\section{第 8 节（测试）}\label{sec:lamet-schemes}

\begin{proposition}[有限 Ioffe 时间窗口内可稳定恢复的自由度]\label{prop:lamet-slepian}
（测试用合成陈述：内容从略。）
\uses{cor:lamet-entire}
\end{proposition}

\begin{example}[截断误差与可恢复模式数的数值]\label{ex:lamet-fourier-numbers}
（测试用合成陈述：内容从略。）
\uses{prop:lamet-truncation, prop:lamet-slepian}
\end{example}

\begin{proposition}[靶质量修正的精确形式 \tagR]\label{prop:lamet-tmc}
（$R_{\mathrm{cond}}$：在所列编号假设之下证明完整；结论的地位不高于其中最弱的假设。）
前提：\cref{asm:lamet-ope}、\cref{thm:lamet-coordinate}、\cref{prop:lattice-renormalizability}。
（测试用合成陈述：内容从略。）
\uses{asm:lamet-ope, thm:lamet-coordinate, prop:lattice-renormalizability, def:lattice-quasi}
\end{proposition}

\begin{example}[$K_N$ 与 $K^\perp_N$（$N\le4$）]\label{ex:lamet-tmc-moments}
（测试用合成陈述：内容从略。）
\uses{prop:lamet-tmc}
\end{example}

\begin{lemma}[Fourier 引理：$z^2\leftrightarrow\partial_x^2$]\label{lem:lamet-fourier}
（测试用合成陈述：内容从略。）
\uses{cor:lamet-entire}
\end{lemma}

\section{第 9 节（测试）}\label{sec:lamet-fourier}

\begin{corollary}[端点二分]\label{cor:lamet-endpoint}
（测试用合成陈述：内容从略。）
\uses{lem:lamet-fourier}
\end{corollary}

\begin{proposition}[扭度 4 修正的重整子模型（大 $\beta_0$） \tagP]\label{prop:lamet-bvz}
（测试用合成陈述：内容从略。）
\uses{lem:lamet-fourier, cor:lamet-endpoint, def:power-borel, prop:power-renormalon-examples}
\end{proposition}
\paragraph{地位与缺口}测试：本陈述的地位由\cref{conj:power-renormalon-ht}、\cref{conj:yzliu-gpd-xi}决定。

\begin{assumption}[扭度展开对 Ioffe 时间的一致性 \tagC]\label{asm:lamet-uniform}
（测试用合成陈述：内容从略。）
地位：测试状态句。
文献：测试。
被使用于：\cref{thm:lamet-hadronic}。
解除途径：测试解除途径。
\uses{thm:lamet-coordinate, prop:lamet-tmc}
\end{assumption}

\begin{assumption}[领头幂次修正的结构（扭度 4） \tagC]\label{asm:lamet-twist4}
（测试用合成陈述：内容从略。）
地位：测试状态句。
文献：测试。
被使用于：\cref{asm:lamet-large-z}、\cref{thm:lamet-hadronic}。
解除途径：测试解除途径。
\uses{prop:lamet-support, cor:lamet-entire}
\end{assumption}

\begin{assumption}[大距离区域的控制 \tagC]\label{asm:lamet-large-z}
（测试用合成陈述：内容从略。）
地位：测试状态句。
文献：测试。
被使用于：\cref{thm:lamet-hadronic}。
解除途径：测试解除途径。
\uses{thm:lamet-coordinate, asm:lamet-twist4}
\end{assumption}

\section{第 10 节（测试）}\label{sec:lamet-power}

\begin{theorem}[强子层次的 LaMET 匹配与幂次修正 \tagR]\label{thm:lamet-hadronic}
（$R_{\mathrm{cond}}$：在所列编号假设之下证明完整；结论的地位不高于其中最弱的假设。）
前提：\cref{asm:lamet-uniform}、\cref{asm:lamet-twist4}、\cref{asm:lamet-large-z}、\cref{thm:lattice-lamet}、\cref{prop:lamet-tmc}。
（测试用合成陈述：内容从略。）
非微扰地位：\cref{asm:intro-wightman}、\cref{asm:intro-composite}。
\uses{asm:lamet-uniform, asm:lamet-twist4, asm:lamet-large-z, thm:lattice-lamet, prop:lamet-tmc, lem:lamet-fourier, cor:lamet-endpoint}
\end{theorem}

\begin{remark}[LaMET 预言的适用域]\label{rem:lamet-domain}
（测试用合成陈述：内容从略。）
% uses: none
\end{remark}

\begin{proposition}[LaMET 与短距离因子化：同一函数的两种展开 \tagR]\label{prop:lamet-sdf}
（$R_{\mathrm{cond}}$：在所列编号假设之下证明完整；结论的地位不高于其中最弱的假设。）
前提：\cref{asm:lamet-ope}、\cref{thm:lamet-coordinate}、\cref{thm:lamet-parton-momentum}、\cref{prop:lattice-renormalizability}。
（测试用合成陈述：内容从略。）
\uses{thm:lamet-coordinate, asm:lamet-ope, thm:lamet-parton-momentum, prop:lattice-renormalizability, prop:lamet-slepian}
\end{proposition}

\begin{remark}[逆问题之争：数学陈述与经验问题]\label{rem:lamet-inverse-problem}
（测试用合成陈述：内容从略。）
% uses: none
\end{remark}

\begin{definition}[准 GPD 与准 DA]\label{def:lamet-quasi-gpd}
（测试用合成陈述：内容从略。）
\uses{def:lattice-quasi, def:excl-gpd, def:excl-da}
\end{definition}

\section{第 11 节（测试）}\label{sec:lamet-sdf}

\begin{theorem}[准 GPD 与准 DA 的匹配 \tagA]\label{thm:lamet-quasi-gpd}
（测试用合成陈述：内容从略。）
\uses{def:lamet-quasi-gpd, asm:lamet-ope, prop:lattice-renormalizability, prop:lamet-support}
\end{theorem}
\begin{proofsketch}测试。\end{proofsketch}
\paragraph{地位与缺口}测试：本陈述的地位由\cref{asm:lamet-ope}、\cref{conj:yzliu-gpd-xi}决定。

\begin{proposition}[其他准分布：重子与重介子 LCDA、$\Delta G$、双部分子分布 \tagP]\label{prop:lamet-other-objects}
（测试用合成陈述：内容从略。）
\uses{def:lattice-quasi, thm:lattice-lamet}
\end{proposition}
\paragraph{地位与缺口}测试：见文献。

\begin{assumption}[订书钉形类空 Wilson 线算符的重整化与 $L\to\infty$ 极限 \tagA]\label{asm:lamet-staple}
（测试用合成陈述：内容从略。）
地位：测试状态句。
文献：测试。
被使用于：\cref{def:lamet-quasi-tmd}、\cref{thm:lamet-quasi-tmd}、\cref{prop:yzliu-qtmd-powercount}。
解除途径：测试解除途径。
\uses{asm:lamet-aux-renorm, def:prelim-wilson}
\end{assumption}
\begin{proofsketch}测试。\end{proofsketch}

\begin{definition}[准 TMD]\label{def:lamet-quasi-tmd}
（测试用合成陈述：内容从略。）
\uses{asm:lamet-staple, def:lattice-quasi, def:prelim-wilson}
\end{definition}

\begin{definition}[离光锥软函数（Ji–Liu–Liu 归一化）]\label{def:lamet-soft-offlc}
（测试用合成陈述：内容从略。）
\uses{def:prelim-wilson, def:tmd-soft, lem:tmd-soft-oneloop}
\end{definition}

\section{第 12 节（测试）}\label{sec:lamet-offforward}

\begin{assumption}[离光锥软函数的大快度结构 \tagA]\label{asm:lamet-soft-rapidity}
（测试用合成陈述：内容从略。）
地位：测试状态句。
文献：测试。
被使用于：\cref{def:lamet-reduced-soft}、\cref{thm:lamet-quasi-tmd}。
解除途径：测试解除途径。
\uses{def:lamet-soft-offlc, prop:tmd-evolution}
\end{assumption}
\begin{proofsketch}测试。\end{proofsketch}

\begin{definition}[约化软函数 $S_r$]\label{def:lamet-reduced-soft}
（测试用合成陈述：内容从略。）
\uses{def:lamet-soft-offlc, asm:lamet-soft-rapidity}
\end{definition}

\begin{assumption}[极限次序 \tagA]\label{asm:lamet-order-limits}
（测试用合成陈述：内容从略。）
地位：测试状态句。
文献：测试。
被使用于：\cref{thm:lamet-quasi-tmd}。
解除途径：测试解除途径。
\uses{def:lamet-quasi-tmd, def:tmd-collins}
\end{assumption}
\begin{proofsketch}测试。\end{proofsketch}

\begin{theorem}[准 TMD 因子化 \tagA]\label{thm:lamet-quasi-tmd}
（测试用合成陈述：内容从略。）
非微扰地位：\cref{asm:intro-bridge}、\cref{asm:intro-composite}。
\uses{def:lamet-quasi-tmd, def:lamet-reduced-soft, def:tmd-collins, prop:tmd-evolution, asm:lamet-staple, asm:lamet-soft-rapidity, asm:lamet-order-limits, lem:tmd-soft-oneloop, prop:tmd-evolution-oneloop, prop:lattice-renormalizability}
\end{theorem}
\begin{proofsketch}测试。\end{proofsketch}
\paragraph{地位与缺口}测试：本陈述的地位由\cref{asm:intro-bridge}决定。

\begin{example}[单圈检验：指数中的因子 $\frac12$]\label{ex:lamet-qtmd-oneloop}
（测试用合成陈述：内容从略。）
\uses{thm:lamet-quasi-tmd}
\end{example}

\section{第 13 节（测试）}\label{sec:lattice-tmd}

\begin{proposition}[格点 Collins–Soper 核 \tagA]\label{prop:lattice-cs-kernel}
（测试用合成陈述：内容从略。）
非微扰地位：\cref{asm:intro-bridge}。
\uses{thm:lamet-quasi-tmd, prop:tmd-evolution}
\end{proposition}
\begin{proofsketch}测试。\end{proofsketch}
\paragraph{地位与缺口}测试：本陈述的地位由\cref{thm:lamet-quasi-tmd}决定。

\begin{proposition}[准 TMD 之比与 TMD 之比 \tagA]\label{prop:lamet-tmd-ratios}
（测试用合成陈述：内容从略。）
\uses{thm:lamet-quasi-tmd}
\end{proposition}
\begin{proofsketch}测试。\end{proofsketch}
\paragraph{地位与缺口}测试：见文献。

\begin{corollary}[LaMET 原理的微扰内容 \tagA]\label{cor:lamet-principle-pt}
（测试用合成陈述：内容从略。）
\uses{thm:lattice-lamet, prop:lamet-gluon-singlet, thm:lamet-quasi-gpd, thm:lamet-quasi-tmd}
\end{corollary}
\begin{proofsketch}测试。\end{proofsketch}
\paragraph{地位与缺口}测试：见文献。

\begin{remark}[批评与回应]\label{rem:lattice-critiques}
（测试用合成陈述：内容从略。）
% uses: none
\end{remark}

\begin{remark}[开放问题]\label{rem:lamet-open}
（测试用合成陈述：内容从略。）
% uses: none
\end{remark}
